\section{Kloosterman emission and two-sided attachment}\label{sec:emitter}
The full-range representation and the critical-range estimate concern different inputs. The former retains the exact positive packets and has a saturation identity; the latter acts on two short arrays. The relevant attachment must therefore be made to the signed, diagonal-removed expression, with explicit recombination and residual costs. This section makes that distinction quantitative.
\subsection{The complete Kloosterman representation}
Let $\eta_b$ be the real partition functions constructed in Section~\ref{sec:boundary}, and define the finite set of physical labels
\[
\mathcal J_x=\{j\in\ZZ:\operatorname{supp}\eta_{j,H}\cap[x/2,x]\ne\varnothing\},
\qquad |\mathcal J_x|=O(x/H+1).
\]
All ordered block pairs and their sums in this section are restricted to $(b,c)\in\mathcal J_x^2$. Set
$a^{(j)}(t)=\eta_b(t)\beta(t)+\ii^j\eta_c(t)w(t)$ for $j=0,1,2,3$.
For fixed $q\in\QQ$ and $v\in\RR$, define the sums on unit residue classes and their centered versions by
\[
F_r^{(j)}=\sum_{t\in I_x,\ t\equiv r(q)}a^{(j)}(t)\e(vt/H),\qquad
F_r^{(j),0}=F_r^{(j)}-\frac1{q-1}\sum_{s\ne0}F_s^{(j)}.
\]
Set the centered vector to zero at the zero residue class. Suppressing the superscript $j$, put
\[
\alpha_m=\frac1q\sum_{r\ne0}F_r^0\e(-mr/q),\quad
\gamma_n=\frac1q\sum_{r\ne0}\overline{F_r^0}\e(-n\bar r/q),\quad 1\le m,n<q,
\]
where $\bar r$ is the multiplicative inverse of $r$ modulo $q$.
\begin{proposition}[Emission identity and saturation]\label{prop:emitter}
For the unnormalized Kloosterman sum $S(m,n;q)=\sum_{r\ne0}\e((mr+n\bar r)/q)$,
\[
\mathcal B=\sum_{m,n=1}^{q-1}\alpha_m\gamma_nS(m,n;q)
=\sum_{r\ne0}|F_r^0|^2
=q\norm\alpha\norm\gamma.
\]
\end{proposition}
\begin{proof}
Centering makes the zero frequency of both Fourier transforms vanish. Expand $S$ and apply finite Fourier inversion to the two inner sums, recovering $F_r^0$ and $\overline{F_r^0}$. Parseval gives
$\norm\alpha^2=\norm\gamma^2=q^{-1}\norm{F^0}^2$, since inversion permutes the unit residue classes.
\end{proof}
This is an exact representation, not a cancellation estimate. The arrays have range $q-1$; centering does not shorten them to length $\sqrt q$. For example, if $F^0=e_1-e_2$, every nonzero frequency of the first array is nonzero. This example tests the interface; it does not identify the actual arithmetic vector with that example.

Appendix~\ref{sec:models} gives the complementary one-sided obstruction: even two first coordinates and a complete centered second array retain the generic norm scale $q$. Neither statement excludes cancellation after signed reassembly of the actual coefficients.

\subsection{Signed recombination before residual estimates}
Write
\[
d_j(q)=\frac{q-2}{q-1}\sum_{t\in I_x,\ q\nmid t}|a^{(j)}(t)|^2,\qquad
E_{bcq}=\int\psi_+(v)\frac14\sum_{j=0}^3\ii^j[\mathcal B_j-d_j(q)]\,dv.
\]
Expanding the four squares and using $\sum_j\ii^j=\sum_j\ii^{2j}=0$ gives
\begin{equation}
\Cx=\sum_{b,c}\sum_{q\in\QQ}qE_{bcq}.\label{eq:polarreturn}
\end{equation}
The phase inside the integral is $\e(v(t-u)/H)$, with the same orientation as the kernel in Section~\ref{sec:operator}.
Subtracting $d_j$ removes the diagonal exactly; the outer factor $q$ must not be counted twice. Each of the four positive packages may saturate even when their signed combination exhibits cancellation. Taking their absolute values separately loses this distinction.

For arrays supported on integer intervals and vanishing at multiples of $q$, write
\[
\mathcal B_q(A_1,A_2)=\sum_{m,n}A_1(m)A_2(n)S(m,n;q).
\]
The unit multiplier is always $1$ here and does not vary with the decomposition index.
When both integer intervals have length at most a constant multiple of $\sqrt q$, the unit multiplier and coprimality hypotheses of the Blomer--Pascadi theorem yield the critical-scale bound
\begin{equation}
|\mathcal B_q(A_1,A_2)|\ll_\epsilon q^{31/32+\epsilon}
\norm{A_1}\norm{A_2}.\label{eq:bp}
\end{equation}
We use the prime-modulus, nonzero-index specialization of their Theorem~1.1~\cite{blomer2026}, not a direct application to the complete emitted arrays.
Cutting each length-$q-1$ array into $k\asymp\sqrt q$ blocks gives
\[
\sum_{A,B}\norm{\alpha_A}\norm{\gamma_B}\le k\norm\alpha\norm\gamma.
\]
Applying \eqref{eq:bp} rectangle by rectangle therefore produces only a $q^{47/32+o(1)}$ norm factor, worse than the full scale in Proposition~\ref{prop:emitter}. This is the cost of that recombination method, not an impossibility theorem for every decomposition.

\begin{condition}[Unproved two-sided attachment with a loss budget]\label{cond:compression}
Fix a loss exponent $\ell\ge0$ and $A_0>2$ before sending $x\to\infty$; let
$\omega_{bc}=(1+|b-c|)^{-A_0}$ for $(b,c)\in\mathcal J_x^2$.
For the already signed, diagonal-removed, integrated quantity $E_{bcq}$, there is a decomposition
\[
E_{bcq}=\sum_\nu\lambda_\nu\mathcal B_q(A_{1,\nu},A_{2,\nu})+R_{bcq},
\]
in which every rectangle satisfies all hypotheses of \eqref{eq:bp}, and
\begin{align*}
\sum_\nu|\lambda_\nu|\norm{A_{1,\nu}}\norm{A_{2,\nu}}
&\ll\omega_{bc}(H/q)x^{\ell+o(1)},\\
|R_{bcq}|&\ll\omega_{bc}Hq^{-1/32}x^{\ell+o(1)}.
\end{align*}
The constants and small-exponent terms are uniform over all original ordered block pairs and the full prime shell, with no unpaid exceptional set.
\end{condition}
If this condition holds, then
\[
\sum_{b,c\in\mathcal J_x}\omega_{bc}
\le |\mathcal J_x|\sum_{h\in\ZZ}(1+|h|)^{-A_0}\ll x/H,
\]
and \eqref{eq:bp}, followed by the outer prime sum, gives
\[
|\Cx|\ll (x/H)H\sum_{q\in\QQ}q^{31/32}x^{\ell+o(1)}
\ll xQ^2x^{-1/96+\ell+o(1)}.
\]
Here $o(1)$ means that an arbitrarily small fixed positive exponent loss can be absorbed at sufficiently large scales: first fix a strict margin, then choose $\epsilon$ in \eqref{eq:bp} sufficiently small.
Thus $\ell<19/2400$ suffices to pay for $\delta>1/400$. This is a conditional exponent ledger. We do not prove Condition~\ref{cond:compression}, nor require every successful route to use this decomposition.
